% problem_id: S001-lemma-quot-Pn
% source_id: S001 (Stacks Project)
% source_file: moduli.tex
% statement_locator: moduli.tex lines 829--839
% proof_locator: moduli.tex lines 841--973
% env: lemma
% id_note: label 在本来源内唯一
% pairing: tight (verified)
% status: complete (statement + author proof verbatim + reason + locators)
%
\begin{lemma}
\label{lemma-quot-Pn}
Let $n \geq 0$, $r \geq 1$, $P \in \mathbf{Q}[t]$.
The algebraic space
$$
X = \Quotfunctor^P_{\mathcal{O}^{\oplus r}_{\mathbf{P}^n_\mathbf{Z}}/
\mathbf{P}^n_\mathbf{Z}/\mathbf{Z}}
$$
parametrizing quotients of $\mathcal{O}_{\mathbf{P}^n_\mathbf{Z}}^{\oplus r}$
with Hilbert polynomial $P$ is proper over $\Spec(\mathbf{Z})$.
\end{lemma}

\begin{proof}
We already know that $X \to \Spec(\mathbf{Z})$ is separated and
locally of finite presentation (Lemma \ref{lemma-quot-s-lfp}).
We also know that $X \to \Spec(\mathbf{Z})$ satisfies the
existence part of the valuative criterion, see
Lemma \ref{lemma-quot-existence-part}.
By the valuative criterion for properness, it suffices to
prove our Quot space is quasi-compact, see
Morphisms of Spaces,
Lemma \ref{spaces-morphisms-lemma-characterize-proper}.
Thus it suffices to find a quasi-compact scheme $T$ and a surjective
morphism $T \to X$. Let $m$ be the integer found in
Varieties, Lemma \ref{varieties-lemma-bound-quotients-free}.
Let
$$
N = r{m + n \choose n} - P(m)
$$
We will write $\mathbf{P}^n$ for
$\mathbf{P}^n_\mathbf{Z} = \text{Proj}(\mathbf{Z}[T_0, \ldots, T_n])$
and unadorned products will mean products over $\Spec(\mathbf{Z})$.
The idea of the proof is to construct a ``universal'' map
$$
\Psi :
\mathcal{O}_{T \times \mathbf{P}^n}(-m)^{\oplus N}
\longrightarrow
\mathcal{O}_{T \times \mathbf{P}^n}^{\oplus r}
$$
over an affine scheme $T$ and show that every point of $X$
corresponds to a cokernel of this in some point of $T$.

\medskip\noindent
Definition of $T$ and $\Psi$. We take $T = \Spec(A)$ where
$$
A = \mathbf{Z}[a_{i, j, E}]
$$
where $i \in \{1, \ldots, r\}$, $j \in \{1, \ldots, N\}$
and $E = (e_0, \ldots, e_n)$ runs through the multi-indices
of total degree $|E| = \sum_{k = 0, \ldots n} e_k = m$.
Then we define $\Psi$ to be the map whose $(i, j)$ matrix
entry is the map
$$
\sum\nolimits_{E = (e_0, \ldots, e_n)}
a_{i, j, E} T_0^{e_0} \ldots T_n^{e_n} :
\mathcal{O}_{T \times \mathbf{P}^n}(-m)
\longrightarrow
\mathcal{O}_{T \times \mathbf{P}^n}
$$
where the sum is over $E$ as above (but $i$ and $j$ are fixed of course).

\medskip\noindent
Consider the quotient $\mathcal{Q} = \Coker(\Psi)$ on $T \times \mathbf{P}^n$.
By More on Morphisms, Lemma
\ref{more-morphisms-lemma-generic-flatness-stratification}
there exists a $t \geq 0$ and closed subschemes
$$
T = T_0 \supset T_1 \supset \ldots \supset T_t = \emptyset
$$
such that the pullback $\mathcal{Q}_p$ of $\mathcal{Q}$ to
$(T_p \setminus T_{p + 1}) \times \mathbf{P}^n$ is flat over
$T_p \setminus T_{p + 1}$. Observe that we
have an exact sequence
$$
\mathcal{O}_{(T_p \setminus T_{p + 1}) \times \mathbf{P}^n}(-m)^{\oplus N}
\to
\mathcal{O}_{(T_p \setminus T_{p + 1}) \times \mathbf{P}^n}^{\oplus r}
\to
\mathcal{Q}_p
\to
0
$$
by pulling back the exact sequence defining $\mathcal{Q} = \Coker(\Psi)$.
Therefore we obtain a morphism
$$
\coprod (T_p \setminus T_{p + 1})
\longrightarrow
\Quotfunctor_{\mathcal{O}^{\oplus r}/\mathbf{P}/\mathbf{Z}}
\supset
\Quotfunctor^P_{\mathcal{O}^{\oplus r}/\mathbf{P}/\mathbf{Z}} = X
$$
Since the left hand side is a Noetherian scheme and the inclusion
on the right hand side is open, it
suffices to show that any point of $X$ is in the image of this morphism.

\medskip\noindent
Let $k$ be a field and let $x \in X(k)$. Then $x$ corresponds to
a surjection $\mathcal{O}_{\mathbf{P}^n_k}^{\oplus r} \to \mathcal{F}$
of coherent $\mathcal{O}_{\mathbf{P}^n_k}$-modules
such that the Hilbert polynomial of $\mathcal{F}$ is $P$.
Consider the short exact sequence
$$
0 \to \mathcal{K} \to
\mathcal{O}_{\mathbf{P}^n_k}^{\oplus r} \to
\mathcal{F} \to 0
$$
By Varieties, Lemma \ref{varieties-lemma-bound-quotients-free}
and our choice of $m$ we see that $\mathcal{K}$ is $m$-regular.
By Varieties, Lemma \ref{varieties-lemma-m-regular-globally-generated}
we see that $\mathcal{K}(m)$ is globally generated.
By Varieties, Lemma \ref{varieties-lemma-m-regular-up}
and the definition of $m$-regularity we see that
$H^i(\mathbf{P}^n_k, \mathcal{K}(m)) = 0$ for $i > 0$.
Hence we see that
$$
\dim_k H^0(\mathbf{P}^n_k, \mathcal{K}(m)) =
\chi(\mathcal{K}(m)) =
\chi(\mathcal{O}_{\mathbf{P}^n_k}(m)^{\oplus r}) -
\chi(\mathcal{F}(m)) = N
$$
by our choice of $N$. This gives a surjection
$$
\mathcal{O}_{\mathbf{P}^n_k}^{\oplus N}
\longrightarrow
\mathcal{K}(m)
$$
Twisting back down and using the short exact sequence above
we see that $\mathcal{F}$ is the cokernel of a map
$$
\Psi_x :
\mathcal{O}_{\mathbf{P}^n_k}(-m)^{\oplus N}
\to
\mathcal{O}_{\mathbf{P}^n_k}^{\oplus r}
$$
There is a unique ring map $\tau : A \to k$ such that the base change
of $\Psi$ by the corresponding morphism $t = \Spec(\tau) : \Spec(k) \to T$
is $\Psi_x$. This is true because the entries of the $N \times r$
matrix defining $\Psi_x$ are homogeneous polynomials
$\sum \lambda_{i, j, E} T_0^{e_0} \ldots T_n^{e_n}$
of degree $m$ in $T_0, \ldots, T_n$ with coefficients
$\lambda_{i, j, E} \in k$ and we can set
$\tau(a_{i, j, E}) = \lambda_{i, j, E}$.
Then $t \in T_p \setminus T_{p + 1}$ for some $p$ and
the image of $t$ under the morphism above is $x$ as desired.
\end{proof}
